\documentclass[11pt,a4paper]{article}
\usepackage[OT1]{fontenc}
\renewcommand{\rmdefault}{cmr}
\renewcommand{\sfdefault}{cmss}
\renewcommand{\ttdefault}{cmtt}
\usepackage{amsmath,amssymb,amsthm,mathtools}
\usepackage[margin=27mm]{geometry}
\usepackage{microtype,booktabs,longtable,array}
\usepackage[hidelinks]{hyperref}
\setlength{\emergencystretch}{2em}
\setlength{\parskip}{3pt plus 1pt}
\setcounter{tocdepth}{1}
\newtheorem{theorem}{Theorem}[section]
\newtheorem{lemma}[theorem]{Lemma}
\newtheorem{proposition}[theorem]{Proposition}
\newtheorem{corollary}[theorem]{Corollary}
\theoremstyle{definition}
\newtheorem{definition}[theorem]{Definition}
\newtheorem{example}[theorem]{Example}
\newtheorem{problem}[theorem]{Open problem}
\theoremstyle{remark}
\newtheorem{remark}[theorem]{Remark}
\newcommand{\N}{\mathbb N}
\newcommand{\NN}{\mathbb N_{>0}}
\newcommand{\Ftwo}{\mathbb F_2}
\newcommand{\lv}{\operatorname{lv}}
\newcommand{\norm}{\operatorname{N}}
\newcommand{\sz}{\operatorname{sz}}
\newcommand{\can}[1]{\mathopen{[\![}#1\mathclose{]\!]}}
\newcommand{\mle}{\preccurlyeq}
\newcommand{\mslt}{\prec}
\newcommand{\PhiF}{\Phi}
\newcommand{\Esrc}{\mathcal E_{\mathrm{src}}}
\newcommand{\headpart}{\operatorname{lim}}
\newcommand{\wt}{\operatorname{wt}}
\newcommand{\Ty}{\operatorname{Ty}}
\newcommand{\FV}{\operatorname{FV}}
\newcommand{\NF}{\operatorname{NF}}
\newcommand{\sgn}{\operatorname{sign}}
\newcommand{\id}{\operatorname{id}}
\newcommand{\BH}{\mathsf{BH}}
\newcommand{\ord}{\operatorname{ord}}
\newcommand{\bodyurl}[1]{\url{#1}}
\numberwithin{equation}{section}

\newcommand{\KK}[2]{K_{#1,#2}}
\newcommand{\EE}[2]{E_{#1,#2}}
\newcommand{\Peak}{\operatorname{Peak}}
\title{Equality and One-Variable Comparison for Positive Collapse Expressions\\
Normal forms and an effective sign algorithm for the standard
Wilken--Weiermann control}
\author{}
\date{28 October 2026\\Working manuscript}
\begin{document}
\maketitle
\begin{abstract}
Fix the standard control $\Phi(t)=F_5(t+100)$, with $F_0(t)=2^t$, in the
Wilken--Weiermann collapse below $\varepsilon_0$.
Consider positive expressions built from natural constants, positive natural
variables, addition and fixed heads $h_\lambda(u)=\psi(\lambda+u)$, where
$\lambda$ is any fixed nonzero limit below $\varepsilon_0$.
We give a complete recursive normal form for universal identities in any
finite number of variables. In one variable we compute both a finite bound
on the exceptional equality set and the eventual sign of every difference.
Exact evaluation of the finite prefix then decides all one-variable Boolean
combinations of equality and order atoms under a single existential or
universal quantifier. The sign algorithm uses exact principal suffix kernels,
retained-mass comparisons and a simultaneous lower-profile separation
certificate. Its termination is proved by a finite reversed-rank ordinal
measure. A source-translation theorem transfers the results to literal
canonical assignments of original G\"odel-T terms with ground free variables.
The identity and one-input results have different scopes; neither supplies
multivariate order comparison or arbitrary existential equality elimination.
\end{abstract}
\tableofcontents

\section{The problem and the positive language}
Wilken and Weiermann's derivation-length analysis assigns ordinal vectors to
terms of G\"odel's T \cite{ww}. Comparing those particular vectors is a
different question from finding some well-founded syntactic ordering of
the calculus. The present paper treats a broad scalar fragment of the
canonical-assignment comparison problem. The companion ordering paper
\cite{wwordering} addresses the syntactic question instead.

Our scalar fragment retains the actual collapse functions. It is not the
elementary language generated by $2^f g$: an algorithm or hardness theorem
for that elementary language does not automatically transfer to expressions
containing $\psi(\lambda+u)$.

Throughout, the norm on Cantor normal forms is additive:
\[
 \norm(0)=0,\qquad
 \norm(\omega^{\alpha_1}+\cdots+\omega^{\alpha_k})
   =\sum_{i=1}^k(1+\norm(\alpha_i)).                 \tag{1.1}
\]
Coefficients are written as repeated summands for this definition. Fix
\[
 F_0(t)=2^t,\quad F_{j+1}(t)=F_j^{\,t+1}(t),\qquad
 f(t)=\Phi(t)=F_5(t+100).
\]
Superscripts on $f$ denote finite iteration, including $f^0(t)=t$.

\begin{definition}[Positive fixed-limit expressions]\label{def:positive}
For a finite list of variables $x_1,\ldots,x_d$ valued in $\NN$, the input
language $\mathcal P_d$ is generated by natural constants, those variables,
addition, and $h_\lambda(A)=\psi(\lambda+A)$, for fixed closed nonzero limits
$\lambda<\varepsilon_0$. Only finitely many heads occur in any input.
Natural coefficient multiplication means repeated addition. There is no
subtraction inside a head, multiplication of two open values, variable
ordinal index, inverse control or additional orbit-height operation.
\end{definition}
A difference $A-B$ is allowed as a signed row to be compared; it is not a
new constructor available inside an input expression. Indices may be
arbitrarily large below $\varepsilon_0$, and nesting is unrestricted.

Three conclusions will be established. Universal equality is decidable in
$\mathcal P_d$ for any finite $d$. In $\mathcal P_1$, a nonidentity has only
effectively bounded finitely many zeros. Finally, in $\mathcal P_1$ its
eventual sign is effective. The third conclusion requires substantially
more than the second: a nonzero integer-valued sequence can change sign
infinitely often without ever being zero.

\subsection{The foundational interface}
We use the effective maximum predecessor, single-predecessor recurrence,
closed evaluation and complete terminal decomposition results of the
companion foundations paper \cite{wwfoundations}, proved there from the
published collapse definition
\cite[\S2]{ww}. Stating the interface also fixes the norm and phase conventions
used by every kernel below.
For a nonzero limit $\eta$ and a natural budget $b$, put
\[
 M(\eta,b)=\max\{\xi<\eta:\norm(\xi)\le b\}.
\]
It is computable in finite Cantor normal form and satisfies
$\norm(M(\eta,b))=b$. For an infinite $\alpha=\eta+n$ the literal next
state of collapse evaluation is
\[
 \alpha\longmapsto M(\eta,f(\norm(\alpha))).         \tag{1.2}
\]
This is strict ordinal descent. If there are $L\ge1$ such steps before the
first finite state and $t=\norm(\alpha)$, then
\[
 \psi(\alpha)=f^L(t)+L.                              \tag{1.3}
\]
The finite case is $\psi(n)=n$. Both the value and the count are computable
by this terminating procedure; no feasible running-time claim is made.

We also use the standard-control estimates and the original controlled
strictness property:
\begin{align}
 f(v)&\ge2^v>v,& f(v+1)&\ge2^{f(v)+1},                \label{eq:gaps}\\
 \alpha<\beta\ \&\ \norm(\alpha)\le f(\norm(\beta))
 &\ \Longrightarrow\ \psi(\alpha)<\psi(\beta).
                                                               \label{eq:strict}
\end{align}
The first line follows from the companion's Elementary control properties
lemma. It is for the nonnegative inputs used here; all elementary
polynomial estimates below are restricted to inputs at least four.
The norm lower bound $\norm(\alpha)\le\psi(\alpha)$ is
\cite[Proposition 2.3]{ww}. These statements do not say that $\psi$ is
unrestrictedly increasing in its ordinal argument.

\begin{lemma}[Sparse peaks]\label{lem:sparse}
For an infinite input of norm $t\ge2$, write $Z=f^L(t)=f(s)$ as in (1.3).
Then $0<L\le s-1$, and
\[
           Z<\psi(\alpha)=Z+L<2Z<f(s+1).             \tag{1.6}
\]
Distinct peaks $f(s)$ and $f(u)$ with $u<s$ satisfy $f(u)\le f(s-1)$.
\end{lemma}
\begin{proof}
Each norm update increases the norm by at least one. Thus
$s=f^{L-1}(t)\ge t+L-1\ge L+1$. Since $f(s)>s$, the count is smaller than
$Z$. The final strict inequality follows from \eqref{eq:gaps}. The last
assertion uses strict increase of the integer map $f$.
\end{proof}

\section{A complete normal form for identities}
\begin{definition}[Collected recursive normal form]\label{def:nf}
A normal form consists of a constant $c\in\N$, a vector
$b\in\N^d$, and a finite dictionary of tags $(\lambda,A)$ with positive
natural coefficients. Each $A$ is itself a normal form with nonempty
variable support. Its value is
\[
       c+\sum_i b_i x_i+\sum_{(\lambda,A)} a_{\lambda,A}h_\lambda(A).
                                                               \tag{2.1}
\]
The dictionary graph is finite and acyclic. The height is zero if no tag
occurs, and otherwise one plus the largest argument height.
\end{definition}
Normalize bottom up. Addition merges dictionaries and coefficient vectors.
For a head, normalize its argument; if that argument is closed, evaluate
the entire call and return a constant. Otherwise keep the tag
$(\lambda,\NF(A))$. Indices are compared as actual canonical Cantor forms.
A fixed finite tail in a presented index is first moved into its argument;
a zero limit part is removed using $\psi(n)=n$.

Identical normal forms have equal values, by induction. Completeness is a
separation theorem and needs more than commutative algebra.

\begin{lemma}[Orbit cancellation at roots]\label{lem:roots}
Call a natural number an $f$-root if it is outside $\operatorname{range}(f)$.
If infinite inputs $\alpha,\beta$ have root norms $t,u$ and equal peaks,
then $\alpha=\beta$.
\end{lemma}
\begin{proof}
From $f^L(t)=f^J(u)$ cancel the smaller number of iterates, using injectivity.
If $L\ne J$, one of the two roots is a positive iterate of the other, a
contradiction. Hence $L=J$ and $t=u$. Formula (1.3) gives equal collapse
values. If the ordinals differed, order them; their common norm satisfies
$t\le f(t)$, so \eqref{eq:strict} would give distinct values.
\end{proof}
This proof is useful because it avoids any assertion that terminal-peak
maps are monotone in arbitrary ordinal inputs.

\begin{theorem}[Identity normal form]\label{thm:identity}
For $E,F\in\mathcal P_d$, the identity $E=F$ holds at every positive
valuation exactly when $\NF(E)=\NF(F)$. If the forms differ, one can
effectively specify a positive countervaluation of the form
\[
           x_i=w_i(f(k)+1),\qquad w_i,k\in\NN.        \tag{2.2}
\]
Consequently universal equality in $\mathcal P_d$ is decidable.
\end{theorem}
\begin{proof}
Let $\mathcal S$ contain the two root normal forms, zero, and all recursively
reachable tag arguments. Let $M$ be the largest raw-variable coefficient
in this finite set, $B=\max(2,M+1)$, and $w_i=B^{i-1}$.
Uniqueness of finite base-$B$ notation shows that different coefficient
vectors in $\mathcal S$ have different weighted slopes $b\cdot w$.
On the ray $x_i=w_i n$, any open positive expression is at least $n$:
this follows by induction from positivity and $h_\lambda(q)\ge q$.

We first make every reachable open atom's initial norm a root on the ray
$n=f(k)+1$. For a tag $(\lambda,A)$ write that norm as
\[
 t= A_0+Dn+\sum_j a_j h_{\lambda_j}(A_j),\quad
 A_0=\norm(\lambda)+c_A\ge2,\quad D=b_A\cdot w,
 \quad C=\sum_j a_j.                                 \tag{2.3}
\]
If $C=0$, openness forces $D\ge1$. For $k\ge\max(4,A_0+2D)$,
\[
          f(k)<t\le(A_0+2D)f(k)<f(k+1),
\]
by \eqref{eq:gaps}. If $C>0$, let $Y=f(s)$ be the largest child peak.
Every child is open, so $Y\ge f(n)\ge2^n$. For
$n\ge\max(4,A_0,D+1,2C+1)$ the affine part is at most $n^2\le Y$.
The largest child contributes more than $Y$, while every child is below
$2Y$ by Lemma~\ref{lem:sparse}. Therefore
\[
                 Y<t<(2C+1)Y<f(s+1).                \tag{2.4}
\]
This too makes $t$ a root. All these bounds are finite syntax data.

Induct on the height of forms in $\mathcal S$. Distinct affine forms are
unequal for all sufficiently large $n$, because their slopes differ or
their constants differ. Suppose distinct lower-height forms have already
been separated. Two distinct tags at the next height cannot have equal
peaks: Lemma~\ref{lem:roots} would identify the ordinals
$\lambda+A(x)$ and $\mu+B(x)$, hence their limit parts and finite tails.
The indices would agree and the different lower-height arguments would
have equal values, contradicting induction.

Now cancel identical tags in the difference of two distinct forms at this
height. If tags remain, write the row on the weighted ray as
\[
                \sum_j c_j h_{\lambda_j}(A_j)+pn+q,  \tag{2.5}
\]
where $c_j\ne0$ and the remaining tags are distinct. Let $Y=f(s)$ be the
largest peak. Its coefficient has absolute value at least one. All other
peaks are at most $U=f(s-1)$, and their full values are below $2U$.
Because the largest tag is open, $s\ge n$, while $U\ge s$. Set
$K=2\sum_j|c_j|+|p|+|q|$. The absolute value of all lower contributions
is at most $KU$. Once $k\ge\max(4,K)$, the leading peak satisfies
$Y\ge2^{U+1}>KU$. Thus the row is nonzero, whichever tag is largest.
No algorithm for determining that largest tag has been assumed.

For explicit effectivity take the maximum of the root bounds above and,
for every pair in $\mathcal S$, the displayed row bound $K$; for an affine
row $pn+q$ with $p\ne0$, also include $|q|+1$. This single computable
$k$ suffices at every induction height. If there are no open variables,
evaluate the two constants instead. Normalization and all finite maxima
terminate, proving the theorem.
\end{proof}

\begin{example}[The exact boundary of positivity]
The collected forms of $h_\omega(x+y)+x$ and $x+h_\omega(y+x)$ agree.
In contrast the two tags in $h_\omega(x)+h_\omega(y)$ remain distinct;
the theorem decides an identity, not whether they happen to coincide at
some valuation. Inner subtraction cannot be added silently. The literal
finite-$\omega$ evaluation gives the exact identity
\[
 h_{2\omega}(n)=h_\omega(h_\omega(n+2)-3)+1
              =f(f(n+4))+2.
\]
The right side's subtraction is forbidden by Definition~\ref{def:positive}.
The free normal form is not asserted for that larger language.
\end{example}

\section{An effective finite-zero bound in one variable}
From now on every open input depends on the same $x\in\NN$.
This hypothesis permits a stronger separation statement along the entire
tail, without selecting the sparse ray (2.2).

\begin{lemma}[Common-band or root separation]\label{lem:oneinput}
For any finite collection of normal forms in $\mathcal P_1$, a computable
$N$ has the following properties for every $x\ge N$.
Every affine open atom input lies in $[x,f(x))$. Every nonlinear atom input
is a root strictly greater than $f(x)$. Equal peaks of two tags force equal
indices and equal argument values.
\end{lemma}
\begin{proof}
Write an atom norm as $A_0+bx+\sum a_jh_{\lambda_j}(A_j)$ as in (2.3).
For an affine input, $b\ge1$ and
$x\le A_0+bx<f(x)$ once $x\ge\max(4,A_0+b+1)$.
For a nonlinear input the largest child peak $Y$ is at least $f(x)$;
for $x\ge\max(4,A_0+b,2C+1)$ the argument proving (2.4) applies
at every such $x$. It is a root greater than $f(x)$.

Two distinct inputs in $[x,f(x))$ cannot lie on the same $f$-orbit,
since the first positive iterate of the smaller is already at least
$f(x)$. Two nonlinear roots are handled by cancellation. In a mixed pair,
orbit cancellation either puts the larger nonlinear root in the range of
$f$, expresses the smaller affine input as an iterate of that larger
input, or identifies their disjoint bands; each is impossible. Thus equal
peaks force equal initial norms and equal iterate counts. Equation (1.3)
and controlled strictness then identify the two ordinals, exactly as in
Lemma~\ref{lem:roots}.
\end{proof}

\begin{theorem}[Effective finite zeros]\label{thm:zeros}
Given $E,F\in\mathcal P_1$, one computes $N\ge1$ such that either their
normal forms agree, or $E(x)\ne F(x)$ for every $x\ge N$. In the latter
case their equality set is the computable finite set obtained by exact
evaluation at $1,\ldots,N-1$.
\end{theorem}
\begin{proof}
Use the same finite set $\mathcal S$ as in Theorem~\ref{thm:identity}.
Take all bounds from Lemma~\ref{lem:oneinput}. For each pair in
$\mathcal S$, cancel its tag dictionaries. If a tag remains, include
$4$ and $2\sum|c_j|+|p|+|q|$ in the maximum; otherwise include
$|q|+1$ when its affine slope $p$ is nonzero. Height induction now
separates distinct argument forms and, by Lemma~\ref{lem:oneinput}, the
peaks of distinct tags. The sparse leading-peak estimate (2.5) separates
the full rows on this entire tail. These finitely many bounds are computed
before the induction and work simultaneously at every height.
\end{proof}

\begin{lemma}[Range cells and signed rows]\label{lem:cells}
The following assertions have effective tail bounds.
For a nonlinear positive expression $E$, let $Z_i$ be its displayed open
tag peaks, and $Y=\max_i Z_i=f(s)$. Then
\[
              Y<E<f(s+1),\qquad E<h_\omega(Y).        \tag{3.1}
\]
For a value $V\in\operatorname{range}(f)$ on that tail,
\[
 V<E\ \Longleftrightarrow\ (\exists i\ Z_i\ge V),\qquad
 V>E\ \Longleftrightarrow\ (\forall i\ Z_i<V).       \tag{3.2}
\]
After cancellation of identical tags in a signed row, its eventual
pointwise sign is the coefficient sign of its greatest remaining peak,
once those peaks have been ordered. If no tags remain, the row is affine.
\end{lemma}
\begin{proof}
Every open child uses the same input $x$, so its peak dominates the entire
fixed affine part on an effective tail. The proof of (2.4), with the
outer index norm omitted, gives $Y<E<(2C+1)Y<f(s+1)$.
Increasing the threshold makes this fixed multiple smaller than
$h_\omega(Y)=f(Y+2)+1$. There is no range value strictly between $f(s)$
and $f(s+1)$, proving (3.2), including impossibility of $V=E$.
Theorem~\ref{thm:zeros}'s height induction separates the remaining tag peaks,
and the estimate in (2.5) proves the signed-row assertion. It does not
identify the greatest peak; that task is supplied below.
\end{proof}

\section{Exact suffix kernels and retained mass}
The sign algorithm uses auxiliary kernels to expose the terminal peak of
a fixed source head. Their indices and masses are closed finite data,
and auxiliary kernels are not new input constructors in $\mathcal P_1$.

\begin{definition}[Principal suffix kernel]\label{def:kernel}
Let $\pi=\omega^\beta\ge\omega$ be a fixed principal, $a\in\N$, and
$t\ge a+\norm(\pi)$. Start at suffix
$\pi+(t-a-\norm(\pi))$, with total norm $t$ and permanent numerical mass
$a$. At a nonfinite suffix $\eta+n$, replace it by
$M(\eta,f(v)-a)$, where $v$ is the current total norm.
Let $\EE\pi a(t)$ count steps to the first finite suffix and put
\[
              \KK\pi a(t)=f^{\EE\pi a(t)}(t).        \tag{4.1}
\]
The finite exit suffix is $\KK\pi a(t)-a$. The value in (4.1) is the
\emph{total} exit norm, and always belongs to $\operatorname{range}(f)$.
\end{definition}
The suffix is strictly descending and its next norm budget is always
genuine. Thus the procedure terminates. A numerical mass need not be the
norm of any source prefix.

\begin{lemma}[Exact prefix splitting]\label{lem:prefix}
Let $\kappa$ be a fixed Cantor prefix and let $\eta$ be a nonzero limit
whose leading monomial is at most the least monomial of $\kappa$;
the zero prefix is also allowed. In particular one may take
$\kappa=\pi k$ and $\operatorname{lead}(\eta)\le\pi$.
For $b\ge c=\norm(\kappa)$,
\[
 M(\kappa+\eta,b)=\kappa+M(\eta,b-c).                \tag{4.2}
\]
A descent above $\kappa$ reaches a state $\kappa+n$ with finite remainder
before passing below $\kappa$. During that segment the remainder's exact
norm control is
\[
                   G_c(v)=f(c+v)-c\ge f(v)>v.         \tag{4.3}
\]
\end{lemma}
\begin{proof}
The ordinal $\kappa$ is eligible below $\kappa+\eta$, so the maximum cannot
have a smaller prefix. Every ordinal in the interval has a unique remainder
$r<\eta$, and compatible Cantor concatenation gives norm $c+\norm(r)$.
Maximizing that remainder proves (4.2). Equal leading boundary exponents
cause coefficients and their norm contributions to add, so they are allowed.
Budgets increase, keeping the prefix affordable. Equation (4.2) forbids
skipping the finite-remainder cut; ordinal well-foundedness supplies it.
Norm saturation proves (4.3). An increasing integer map preserves integer
differences, giving its inequality.
\end{proof}

\begin{lemma}[Ordered controls and counts]\label{lem:counts}
For increasing integer controls $g\le h$, both strictly exceeding their
input, let $L_g,L_h$ be the respective norm-maximal descent counts.
Then
\[
 \alpha\le\beta,quad\norm(\alpha)\le\norm(\beta)
       \quad\Longrightarrow\quad L_g(\alpha)\le L_h(\beta). \tag{4.4}
\]
\end{lemma}
\begin{proof}
Induct on $\beta$. Finite $\alpha$ has zero count. Otherwise write the
limit parts of $\alpha,\beta$ as $\lambda,\mu$. If $\lambda=\mu$, their
next states $p,q$ are ordered by monotonicity in the budget, and
$\norm(p)=g(\norm\alpha)\le h(\norm\beta)=\norm(q)$.
If $\lambda<\mu$, then $\alpha<\mu$ and its norm is less than the next
$h$-budget. Hence $\alpha$ is eligible below $\mu$, so
$p<\alpha\le q$; the same norm inequality holds. Apply induction at
$q<\beta$ and add the first step.
\end{proof}
This also proves fixed-head monotonicity: increasing the finite tail of
$\rho+n$ increases its norm and does not decrease its count, so its full
$\psi$ value is strictly increasing. Fixed-mass kernels strictly increase
their genuine input by the same argument. Decreasing their retained mass
at fixed total input weakly increases the suffix and its budgets, so cannot
decrease the exit norm.

\begin{lemma}[Exact peak word]\label{lem:peakword}
For each fixed limit $\lambda$, one computes a finite word $W_\lambda$ of
principal kernels such that the peak of $h_\lambda(A)$ is exactly
\[
             Z_{\lambda,A}=W_\lambda(\norm(\lambda)+A). \tag{4.5}
\]
The omitted count is positive: $h_\lambda(A)=Z_{\lambda,A}+D_{\lambda,A}$.
\end{lemma}
\begin{proof}
Expand $\lambda$ into coefficient-one principal Cantor blocks. Delete them
from right to left. For each block $\pi$, the remaining earlier prefix has
a fixed norm $a$. Lemma~\ref{lem:prefix} says that this segment changes the
current total norm by exactly $\KK\pi a$. The last segment has prefix
zero. Concatenating the finitely many segments gives the exact word.
Every removed nonfinite segment contributes its literal count to $D$;
no count is set to zero in (4.5).
\end{proof}

\section{Paying an entire lower positive profile}
A lower profile below a principal $L>\omega$ means a fixed expression
$P(z)$ formed from constants, $z$, addition and heads $h_\rho$ with
$\rho<L$. Its heads may themselves be nonprincipal limits.
The following reserve argument is needed to compare unequal kernel inputs
and to handle a limit exponent. Merely knowing that a larger kernel exceeds
a fixed number of $f$-iterates would be insufficient.

\begin{lemma}[One-stage reserve consumer]\label{lem:stage}
Let $\rho<\pi m$ be a fixed nonzero limit, $m\ge1$, and begin an ordinary
descent at $A=\pi(k+m)+n$, of norm $X\ge x\ge\norm(\rho)$.
At its first cut $\pi k+n'$ with finite remainder,
\[
        n'>J_\rho(x):=f^{D_\rho(x)}(x)+D_\rho(x),    \tag{5.1}
\]
where $D_\rho(x)$ is the count from
$\rho+(x-\norm(\rho))$.
\end{lemma}
\begin{proof}
Put $c=k\norm(\pi)$ and $\alpha=\rho+(x-\norm\rho)$.
The ordinal $\pi k+\alpha$ is below the initial limit part and has norm
$c+x\le2X<f(X)$, hence is eligible for the first predecessor. Its
remainder is at least $\alpha$ and has norm $f(X)-c\ge x$.
Lemmas~\ref{lem:prefix} and \ref{lem:counts} show that the following
prefix segment takes at least $d=D_\rho(x)$ further steps. Its exit total
norm $Y$ is at least $f^{d+1}(X)$. As $c\le X$ and every increasing
integer iterate preserves differences,
\[
 n'=Y-c\ge f^{d+1}(X)-X\ge f^{d+1}(x)-x.
\]
Write $W=f^d(x)\ge x+d$. Then
$f(W)-x>2W-x\ge W+d=J_\rho(x)$, as required.
Notice that the prefix mass was paid using $c\le X$; it need not be
bounded by the smaller test input $x$.
\end{proof}

\begin{lemma}[Effective profile reserve]\label{lem:reserve}
From a lower profile $P$ below $L>\omega$ one computes a nonempty word
\[
 F=h_{\rho_l}\circ\cdots\circ h_{\rho_1},\quad
 F(z)>P(z)\ (z\ge z_0),\quad
 R=\sum_j(\norm(\rho_j)+1),\quad
 H=\max_j\norm(\operatorname{lead}(\rho_j)),          \tag{5.2}
\]
with every $\rho_j<L$. If $\pi\ge\omega$ is principal and at least every
displayed leading monomial, a descent from $\pi R+n$, $n\ge z\ge z_0$,
has terminal norm greater than $P(z)$. The same conclusion holds with any
additional permanent numerical mass $c$.
\end{lemma}
\begin{proof}
Choose $z_0\ge4$, all constants and every weighted addition fan-in in $P$.
Process its expression graph in topological order, maintaining a common
upper bound $u$ on all node values already visited. An addition with fan-in
$b$ is bounded by $bu\le u^2\le2^u<h_\omega(u)$; append $h_\omega$ to the
running bound. For a head node append that head, using its monotonicity.
Every update increases the running bound. Append a final $h_\omega$ to
make the bound strict and the word nonempty. Since $L>\omega$, all added
heads are still below $L$.

For a head $\rho_j$ take $m_j=\norm(\rho_j)+1$ blocks. Its leading
coefficient is at most its norm, so $\rho_j<\pi m_j$. Moreover the
current norm at $\pi(k+m_j)+n$ is at least $\norm(\rho_j)+n$.
Lemma~\ref{lem:stage}, with $x=\norm(\rho_j)+n$, therefore makes the
finite tail at the next cut exceed $h_{\rho_j}(n)$. Repeat through the
word, using exactly $R$ blocks. The final finite norm exceeds $F(z)>P(z)$.
For permanent mass $c$, the exact suffix control is $G_c\ge f$.
Lemma~\ref{lem:counts} orders the counts and the iterated norms, so the
terminal suffix norm is at least the unprefixed one; adding $c$ only
increases the total norm.
\end{proof}

\begin{lemma}[Direct kernel growth]\label{lem:direct}
For fixed $L>\omega$ principal and fixed lower profile $P$, one computes
$T$ such that, for $t\ge T$ and every genuine mass $0\le c\le t-\norm(L)$,
\[
                         \KK Lc(t)>P(t).             \tag{5.3}
\]
Also, $\KK\pi c(t)\le h_\pi(t)$ on its genuine domain.
\end{lemma}
\begin{proof}
Choose one fixed principal $\omega\le\pi<L$ dominating the finite list of
leading monomials in (5.2). For
$t\ge\max(z_0,R\norm\pi,4,\norm L)$, the ordinal $S=\pi R+t<L$ has
\[
             c+\norm(S)\le R\norm\pi+2t\le3t<f(t).
\]
It is eligible below the actual first suffix $M(L,f(t)-c)$. Count
comparison under the same induced control transfers the reserve bound
from $S$ to that maximum. This proves (5.3) uniformly in the masses.
For the final assertion first lower the retained mass to zero, then
increase the seed from $t$ to $\norm\pi+t$. Mass monotonicity and
fixed-mass input monotonicity bound the kernel peak by the full head value.
\end{proof}

\section{Adjacent inputs and limit-exponent mass separation}
\begin{theorem}[Cross-mass strict-input separation]\label{thm:adjacent}
For a fixed principal $L>\omega$ and fixed lower profile $P$ with reserve
data (5.2), put
\[
              T=\max(4,R+5,z_0,\norm L,H).            \tag{6.1}
\]
For $t\ge T$, $0\le c\le t+1-\norm L$, and
$0\le e\le t-\norm L$,
\[
       \KK Lc(t+1)>P(\KK Le(t)+\EE Le(t)).            \tag{6.2}
\]
There is no ordering requirement on $c,e$. More generally the same
conclusion holds with any $u>t$ and any mass genuine at $u$ on the left.
\end{theorem}
\begin{proof}
Write $B=f(t)$, $B'=f(t+1)$, $b=M(L,B-e)$,
$\pi=\operatorname{lead}(b)$ and $w=\norm\pi$. Since
$B-e\ge f(t)-t+\norm L\ge t+\norm L\ge H$, every leading monomial
needed by (5.2) is eligible below $L$ in this budget. Also $B-e\ge2$,
so $\pi\ge\omega$. Thus $\pi<L$, $w\le B$, and $\pi$ dominates
every required leading monomial.

Put $p=c+Rw$ and
\[
 k=B'-B-(p-e),\qquad A=\pi R+b+k.
\]
Here $p-e\le t+1+RB\le(R+1)B$ and
$B'\ge2^{B+1}>(R+2)B$, so $k\ge0$. The sum is compatible CNF,
has norm $B'-c$, and lies below the principal $L$. Hence
$A\le Q=M(L,B'-c)$ and both have the same norm. The latter is the
upper process's actual first suffix.

Follow the packed suffix $A$ retaining its first $R$ blocks while deleting
the remainder $b+k$. At the reserve cut it has finite tail
\[
               n=V-p,\qquad V=f^q(B').              \tag{6.3}
\]
The lower process, after its first step, starts at $b$ with norm $B$
and retained mass $e$. Let its remaining count be $\ell$, and set
$D=1+\ell=\EE Le(t)$ and $Z=f^\ell(B)=\KK Le(t)$.
The packed remainder initially exceeds $b$. The initial global gap pays
the mass difference $p-e$ and, as increasing integer iterates preserve
differences, at every simultaneous later step
$f^i(B')-p\ge f^i(B)-e$. Monotonicity of $M$ in its limit bound and
budget keeps the packed remainder above the lower remainder; it cannot
become finite earlier. Hence $q\ge\ell$.

If $s=f^{D-1}(t)$, then
\[
 V\ge f^\ell(B')=f^D(t+1)\ge f(s+1)\ge2^{Z+1}.
\]
Furthermore $D<Z$ and $p\le(R+1)Z$. Since $Z\ge t\ge R+5$,
$p+Z+D<(R+3)Z<Z^2<2^{Z+1}\le V$. Thus $n>Z+D$ after the full
retained mass has been subtracted. Lemma~\ref{lem:reserve} gives terminal
total norm greater than $P(Z+D)$ for the packed process. Count comparison
from $A\le Q$ at their equal norms, under their common control $G_c$,
transfers that bound to the actual upper kernel.

For arbitrary $u>t$, use the worst-mass inequalities
\[
 \KK Lc(u)\ge \KK L{u-\norm L}(u)
          \ge \KK L{t+1-\norm L}(t+1).              \tag{6.4}
\]
The first is mass monotonicity. For the second both processes start at
the identical suffix $L$, while their simultaneous suffix budgets are
$f^i(u)-u+\norm L\ge f^i(t+1)-(t+1)+\norm L$.
Their counts and terminal total norms are ordered. Apply the adjacent
case with the displayed genuine upper mass.
\end{proof}

At $L=\omega$ the kernel is simply $f(t)$, independently of mass. There
are no nonzero lower limit heads, so $P(z)=az+b$. Taking
$T_\omega=\max(4,a+b+2)$ gives
$P(f(t)+1)<f(t+1)$ by \eqref{eq:gaps}. This is the separate base case of
Theorem~\ref{thm:adjacent}.

\begin{theorem}[Limit-exponent mass separation]\label{thm:limitmass}
Let $L=\omega^\beta$ with $\beta$ a nonzero limit. For a fixed lower
profile $P$ and data (5.2), put
\[
           T=\max(4,R+4,z_0,\norm L+1,H).            \tag{6.5}
\]
For $t\ge T$ and $0\le c<e\le t-\norm L$,
\[
                \KK Lc(t)>P(\KK Le(t)+\EE Le(t)).    \tag{6.6}
\]
Consequently the lexicographically larger pair $(\text{input},-\text{mass})$
gives domination against every fixed lower profile, on an effective
genuine tail. Equal pairs give equal values.
\end{theorem}
\begin{proof}
Put $B=f(t)$ and $C=f(B)$. Because $\beta$ is a limit, the maximum
affordable exponent $\delta_a=M(\beta,B-a-1)$ saturates its norm budget.
Its monomial $\pi_a=\omega^{\delta_a}$ already costs $B-a$, so the first
suffix below $L$ is exactly $\pi_a$, with coefficient one and no remainder.
Since $c<e$, saturation gives $\delta_c>\delta_e$ and $\pi_c>\pi_e$.
Put $w=\norm\pi_e=B-e$. The inequality
$B-e\ge f(t)-t+\norm L\ge t+\norm L\ge H$ shows that $\pi_e$ is
infinite and dominates every required leading monomial.

At the second upper step pack
\[
       k=C-c-(R+1)w,\qquad A=\pi_e(R+1)+k.
\]
Since $C>(R+1)B$ and $c+(R+1)w\le(R+1)B$, the tail $k$ is nonnegative.
The ordinal $A$ is below $\pi_c$ and has norm $C-c$, hence is below the
actual second upper suffix $Q=M(\pi_c,C-c)$ at that same norm.

In the packed process hold $R$ blocks while deleting $\pi_e+k$.
The retained mass is $p=c+Rw$. Compare this remainder with the lower
process after its first step, namely $\pi_e$ with total norm $B$ and
retained mass $e$. The global difference $C-B$ exceeds $p-e\le RB$,
so synchronized budgets keep the packed remainder above the lower one.
If the latter needs $\ell$ further steps, the reserve cut has norm
$V\ge f^\ell(C)=f(Z)$, where $Z=\KK Le(t)=f^\ell(B)$ and
$D=\EE Le(t)=1+\ell<Z$. Since $p\le(R+1)Z$ and $Z\ge R+4$,
$f(Z)>(R+3)Z>p+Z+D$. Its finite tail after all retained mass is paid
is therefore greater than $Z+D$.

Consume the reserve by Lemma~\ref{lem:reserve}, still with permanent mass
$c$. Transfer from $A\le Q$ at equal suffix norms using
Lemma~\ref{lem:counts}. This yields (6.6). Every estimate uses only the
displayed genuine domains, not a fixed bound on $c,e$, so it is uniform
over the stated numerical masses. Combine with
Theorem~\ref{thm:adjacent} for unequal inputs to obtain the lexicographic
assertion. A profile-dependent threshold is intended; there is no single
threshold valid for all profiles at once.
\end{proof}
The limit-exponent hypothesis is essential. For example, the successor
identity proved next gives
$\KK{\omega^2}0(t)=f(\KK{\omega^2}2(t))$, while the lower profile
$h_\omega(z)=f(z+2)+1$ exceeds $f(z)$. Thus arbitrary lower-profile
absorption at the same input is false for successor exponents.

\section{Successor cuts and all same-input mass classes}
\begin{lemma}[Exact successor cut and count cocycle]\label{lem:successor}
Let $\theta=\omega^\gamma\ge\omega$, $L=\theta\omega$, and
$w=\norm\theta$. On the common genuine domain $t\ge c+w+\norm L$,
\begin{align}
 \KK Lc(t)&=\KK\theta c(\KK L{c+w}(t)),              \label{eq:cut}\\
 \EE Lc(t)&=\EE L{c+w}(t)
        +\EE\theta c(\KK L{c+w}(t)).                \label{eq:cocycle}
\end{align}
\end{lemma}
\begin{proof}
For $b=wq+r$, $0\le r<w$, the literal maximum is
\[
                   M(L,b)=\theta q+M(\theta,r).      \tag{7.3}
\]
The largest monomial below $L$ is $\theta$, costing $w$; greedily maximize
its coefficient, then the remainder below $\theta$. This also covers
$q=0$. Lowering the first budget from $f(t)-c$ to $f(t)-c-w$ removes
exactly one $\theta$ block and leaves the same remainder. Retaining that
one block in the $c$-process makes its mass during deletion of the other
remainder exactly $c+w$. Both processes have total norm $f(t)$ after
their first step. Prefix splitting therefore gives identical subsequent
norms until the latter process's finite cut. The remaining block is
exactly $\KK\theta c$, and the counts add. The stronger common domain
ensures that the inner shifted-mass kernel exists.
\end{proof}

\begin{lemma}[Small principal kernels]\label{lem:small}
On their genuine domains,
\[
 \KK\omega c(t)=f(t),\quad\EE\omega c(t)=1,\qquad
 \EE{\omega^2}c(t)=f(t)/2-\lceil c/2\rceil+1.         \tag{7.4}
\]
Thus $\omega^2$ kernels at the same input have exactly the mass classes
$\lceil c/2\rceil$, ordered in reverse. The control has even values.
\end{lemma}
\begin{proof}
Below $\omega$ the first predecessor is finite. Below $\omega^2$ it is
$\omega q+r$ with $q=\lfloor(f(t)-c)/2\rfloor$ and $0\le r<2$.
The first step is followed by exactly one step for each $\omega$ block,
because each such block's deletion uses a finite suffix maximum. Thus
the count is $q+1$, giving (7.4). The definition of the iterated
exponential hierarchy makes each $f(t)$ even.
\end{proof}

\begin{theorem}[Strict mass order above degree two]\label{thm:mass}
For every fixed principal $L\ge\omega^3$, a computable $T_L$ satisfies
\[
 t\ge T_L,quad0\le c<e\le t-\norm L
                 \quad\Longrightarrow\quad\KK Lc(t)>\KK Le(t). \tag{7.5}
\]
For a nonzero limit exponent take $T_L=\max(7,\norm L+1)$.
For $L=\theta\omega$, $\theta\ge\omega^2$, $w=\norm\theta$, take
\[
 T_L=\max(4,\norm L,2w,\max(8,w),S_\theta),          \tag{7.6}
\]
where $S_\theta=0$ at $\theta=\omega^2$ and otherwise $S_\theta=T_\theta$.
\end{theorem}
\begin{proof}
The limit case is Theorem~\ref{thm:limitmass} with profile $P(z)=z$.
Its compiler uses just one final $h_\omega$, giving $R=3,H=2,z_0=4$.
For the successor case let $B=f(t)$, and write the first suffixes from
(7.3) as $\theta q_a+r_a$, $a=c,e$. Here $q_c\ge q_e=q$ and
$r_a=M(\theta,(B-a)\bmod w)$. The domain and (7.6) imply
$B-e\ge f(t)-t+\norm L\ge2w$, so $q\ge2$.

Hold $q$ common $\theta$ blocks while deleting residuals
$\theta(q_c-q)+r_c$ and $r_e$. The first residual is at least the
second. Their permanent masses are $c+wq$ and $e+wq$; at equal global
norms their suffix budgets preserve this order. The finite-cut entry
norms therefore satisfy $U_c\ge U_e\ge B$.

With $k$ common blocks remaining, the child masses are
$a_k=c+w(k-1)<b_k=e+w(k-1)$. They are genuine even at the smaller
entry norm $V\ge B$, because
\[
 b_k\le e+w(q-1)=B-((B-e)\bmod w)-w\le B-w\le V-w. \tag{7.7}
\]
If the entry norms are unequal, Theorem~\ref{thm:adjacent} with identity
profile preserves strict order, uniformly over both masses; its threshold
is $\max(8,w)$. If they agree and $\theta\ge\omega^3$, the inductive
strict-mass theorem makes the first child outputs strictly ordered, and
every subsequent block preserves it.

For $\theta=\omega^2$, Lemma~\ref{lem:small} gives weak order, with a
tie precisely when $b_k=a_k+1$ and $a_k$ is odd. If the first block ties,
the next pair of masses subtracts $w=3$ from both; the smaller mass is
then even, so their ceilings differ and the second outputs are strictly
ordered. There are at least two blocks, and (7.7) pays their domains.
All further blocks preserve the strict order. This proves the base and
the successor induction. Removing the finite terminal successor tail of
the exponent reaches either this base or a limit exponent, so the
threshold recursion terminates.
\end{proof}

\section{The complete one-variable comparison algorithm}
We now supply the order decision that the normal-form proof deliberately
did not assume. There are two sorts in the auxiliary comparison language.
The source sort $P$ is $\mathcal P_1$. A range expression $R$ is a
nonempty finite word of kernels applied to one open $P$ expression.
Every argument of a kernel is either $P$ or $R$; a sum of range expressions
is not introduced into a kernel argument. Closed words are evaluated and
folded into constants. Every open expression of either sort is at least
$x$ on a computable tail. All generated masses and indices are fixed
natural or ordinal data on each branch.

\begin{theorem}[Effective eventual signs]\label{thm:sign}
Given two expressions in these sorts, one computes a sign
$\sigma\in\{-1,0,1\}$ and $N\ge1$ such that their difference has sign
$\sigma$ for every $x\ge N$. In particular this holds for every pair in
$\mathcal P_1$.
\end{theorem}
The proof is a simultaneous comparison and separation-certificate algorithm.
We specify its rules before proving their common termination measure.
Every rule retains the thresholds of all its children and the finitely
many kernel domains it uses. Since every open argument is at least $x$
eventually, each fixed genuine-domain obligation is met by increasing
$N$ to a computable constant.

\subsection{The comparison rules}
For $P/P$, form the signed normal-form row. An affine row is decided
directly. If tags remain, replace each relevant tag by its exact peak
word (4.5), compare those $R$ expressions recursively and choose the
greatest peak. Distinct tags cannot have equal peaks beyond the bound of
Theorem~\ref{thm:zeros}. Return its coefficient sign, using
Lemma~\ref{lem:cells}. Identical forms return exact equality.

For $R/P$, a constant or nonconstant affine $P$ is eventually smaller:
the open $R$ exceeds $f(x)$, which dominates every fixed affine expression.
For a nonlinear $P$ compare $R$ with its displayed peaks (4.5), and use
the exact range-cell test (3.2). The $P/R$ rule is symmetric.

For $R/R$ with equal outer principal, compare the two literal arguments
recursively. Strictly ordered arguments give the same strict kernel order
by Theorem~\ref{thm:adjacent}, or by monotonicity of $f$ at $\omega$.
If the arguments agree eventually, use the following same-input key:
mass is irrelevant at $\omega$; compare $\lceil c/2\rceil$ in reverse at
$\omega^2$; compare masses in reverse above $\omega^2$ by
Theorem~\ref{thm:mass}. Equal keys give equal outer values.

For different outer principals $H>M$, if $H$ has a successor exponent,
use the exact cut \eqref{eq:cut} on the $H$ side until the new outer
principal is $M$ or has a limit exponent still above $M$.
This cannot jump below $M$: every principal below $\theta\omega$ is at
most $\theta$. It takes only the finite terminal successor part of the
exponent. Combine these rewrites immediately with the equal-principal
rule or the following mixed-limit rule; do not recursively call the
unmodified comparison at an intermediate, larger syntax state.

In the mixed-limit case write
\[
                  X=\KK Lc(A),\qquad Y=\KK Me(B),\quad M<L,
                                                               \tag{8.1}
\]
where $L$ has a nonzero limit exponent. Recursively compare $X$ with $B$.
If $X\le B$, then $Y>B\ge X$. If $X>B$, the separation certificate
below, with lower profile $h_M$, gives
$X>h_M(B)\ge Y$, by Lemma~\ref{lem:direct}. There is no equality branch
in this mixed-limit case.

\subsection{The simultaneous limit-cut certificate}
Whenever one original side of a comparison is a bare limit-exponent
kernel $X=\KK Lc(A)$ and its result is $X>S$, the algorithm also
constructs, for every requested fixed lower profile $P$ below $L$, a
computable modulus $N_{X,S}(P)$ such that
\[
                   X(x)>P(S(x))\quad(x\ge N_{X,S}(P)). \tag{8.2}
\]
This concerns generated $P/R$ expressions $S$, not arbitrary integers
below $X$: the assertion for $S=X-1$ would be false, and subtraction
does not occur in the generated language.

Here is the certificate's complete case analysis, following the same
strictly smaller children as comparison.
\begin{enumerate}
\item If $S=ax+b$, use direct growth (5.3) for the fixed profile
$z\mapsto P(az+b)$ at the actual argument $A\ge x$. Positivity of $P$
makes this imply (8.2).
\item If $S$ is nonlinear positive, $S<X$ forces every child peak $Z_i<X$
by the range-cell rule. Apply each smaller child's certificate to
$P\circ h_\omega$ and use $S<h_\omega(\max_i Z_i)$. Take the finite
maximum of their moduli; no permanent choice of maximizing child is needed.
\item If $S$ has the same outer principal $L$, strict order means either
a greater input for $X$, or equal inputs and smaller mass. The
lexicographic conclusion of Theorems~\ref{thm:adjacent} and
\ref{thm:limitmass} applies directly to $P$, at the genuine integer inputs.
\item If $S=\KK Me(B)$ has smaller outer principal, its comparison used
the smaller child $X>B$. Apply that child's certificate to
$P\circ h_M$ and use $S\le h_M(B)$. This handles any finite nesting of
lower contexts by recursive composition, not by an unproved absorption step.
\item If $S$ has larger limit outer principal $H>L$, its comparison with
$X$ compares $S$ against $A$. The only branch giving $X>S$ has $S\le A$;
direct growth for $\KK Lc$ gives $X>P(A)\ge P(S)$. In the other branch
the larger kernel's child certificate would force
$S>h_L(A)\ge X$, so cannot yield $X>S$.
\item A larger successor outer principal on $S$ is unfolded exactly until
it reaches $L$ or a larger limit exponent; use the corresponding preceding
case. The rewrites are exact, including their changed masses and domains.
\end{enumerate}
Every requested profile remains a fixed positive expression with heads
below the relevant limit $L$: only affine maps, fixed lower heads and
$h_\omega$ are composed. Thresholds are profile dependent, and are computed
by following these finitely many smaller children.

\subsection{A common well-founded measure}
Expand sharing when counting occurrences. Let $h$ be the number of
original source $h_\lambda$ nodes across the two compared expressions,
including their argument nodes. Let $\mathcal Q$ be the finite collection
of principal indices in the Cantor decompositions of all original source
heads and initial auxiliary kernels. Close it under finite successor-exponent
chains, stopping at exponent one or a limit exponent. No rule introduces
a principal outside this finite set. List it increasingly as
$\pi_0<\cdots<\pi_{m-1}$ and assign a kernel occurrence of index $\pi_i$
weight $\omega^{m-1-i}$. Let $W$ be their natural sum across the pair.
Use the lexicographic measure
\[
                             (h,W).                 \tag{8.3}
\]

Replacing an original tag by its peak removes that displayed $h$ node,
retaining its argument only once in each child comparison. In $P/P$ and
$R/P$ children, other siblings are not duplicated. Thus $h$ decreases.
Equal-principal comparison removes both outer kernels, decreasing $W$.
Mixed-limit comparison removes the lower outer kernel, also decreasing $W$.
A combined mixed-successor step removes the opposite lowest-rank $M$
occurrence and adds only finitely many occurrences of principals strictly
greater than $M$, retaining the original higher occurrence internally.
Their total weight is strictly smaller than one $M$ weight. Hence $W$
decreases. Affine or closed leaves stop.

Since $h$ is natural and $W<\omega^m$, every branch terminates. Branching
is finite, so the comparison tree is finite. The certificate uses that
same tree: each invocation in its case analysis belongs to a strictly
smaller comparison child or is a direct supplier application. Thus no
unchanged comparison state calls its own certificate. Mutual induction
on (8.3) proves both comparison correctness and (8.2), establishing
Theorem~\ref{thm:sign}.

\subsection{Closure and exact finite exceptions}
The two sorts are closed under every recursive child. Peak replacement
introduces a finite kernel word on a positive argument plus its fixed
norm constant. Successor expansion introduces only a longer kernel word.
No signed inner offset, exposed count, division, inverse iterate or
variable ordinal index occurs. Increasing a mass changes fixed domain
constants without changing the finite rank set or measure.

\begin{example}[A nested comparison with the counts retained]
Consider $E=h_{2\omega}(x)$ and $F=h_\omega(h_\omega(x))$.
Deleting the two $\omega$ blocks gives
\[
 E=f(f(x+4))+2,\qquad F=f(f(x+2)+3)+1.
\]
Thus their peaks are respectively $f(f(x+4))$ and
$f(f(x+2)+3)$. The equal-outer-$\omega$ rule compares the arguments
$f(x+4)$ and $f(x+2)+3$. The latter positive nonlinear expression lies
strictly between its peak $f(x+2)$ and the next range value on the
computed tail. Comparing $f(x+4)$ with that peak reduces to the affine
comparison $x+4>x+2$. Hence $f(x+4)>f(x+2)+3$, and the original signed-row
rule returns $E>F$ on the effective tail. The offset $3$ has been paid
by a range-cell test, rather than deleted as negligible without proof.
The final procedure evaluates the finitely many earlier inputs exactly.
\end{example}

\begin{example}[Why the mass key changes at degree three]
At the same genuine input, the masses $1$ and $2$ give equal
$\omega^2$ kernels because their ceilings after division by two agree.
For a parent $\omega^3$ comparison with mass difference one, the first
common $\omega^2$ block can likewise tie when its retained child masses
are an odd $a$ and $a+1$. The next block reduces both child masses by
three; its ceiling classes are then distinct. The guaranteed two common
blocks ensure that those reduced masses remain nonnegative.
This is the two-block mechanism in
Theorem~\ref{thm:mass}, and explains why the algorithm uses neither strict
mass order at degree two nor a ceiling equivalence at all higher degrees.
\end{example}

\begin{corollary}[One-input quantified comparison]\label{cor:boolean}
For any finite Boolean combination of equality, weak-order and strict-order
atoms between expressions of $\mathcal P_1$, each of its existential and
universal closures over $x\in\NN$ is decidable.
\end{corollary}
\begin{proof}
Apply Theorem~\ref{thm:sign} to every atom and take the maximum tail bound.
The Boolean truth value is constant beyond that bound. Evaluate the
original positive input expressions exactly at every earlier positive
integer. For an existential statement, a true prefix instance or a true
tail suffices; for a universal statement require all prefix instances and
the tail. Evaluate original expressions, rather than auxiliary kernels
before their raised genuine domains.
\end{proof}

\section{Consequences for canonical assignments and exact limits}
The Actual-source forward translation theorem in the companion foundations
paper \cite{wwfoundations} proves the following assertion for original
undecorated G\"odel-T terms. If all free
term variables have ground type, every positive canonical coordinate is
a closed computable ordinal, and coordinate zero is effectively translated,
valuation preservingly, into $\mathcal P_d$. The source may have arbitrary
temporary higher-type binders, recursors, non-normal syntax and output
type. The result is a generating invariant of the actual source translation;
it is not an assertion that every ambient ordinal-expression grammar has
closed positive coordinates. In particular, it does not include the
decorated auxiliary $R_t$ syntax of the extended rewrite system.

\begin{corollary}[Ground-open source equality]\label{cor:sourceeq}
Universal equality of literal canonical assignment vectors is decidable
for any original source pair whose free names all have ground type.
\end{corollary}
\begin{proof}
Construct their canonical assignments by the literal syntax rules. Compare
the closed positive coordinates exactly and compare their zero-coordinate
normal forms by Theorem~\ref{thm:identity}. A positive mismatch refutes
vector equality. Equal positive coordinates and equal normal forms prove
it. If the normal forms differ, (2.2) supplies an admissible positive
ground-coordinate countervaluation. This is coordinate substitution, not
replacing a free variable by a source numeral and recomputing its assignment.
\end{proof}

For different vector levels, pad with zeros through their maximum level,
as in the original assignment convention. If $A,B$ have only one shared
ground free name $X$, write their residuals as $E_A(x),E_B(x)$ and put
\[
 U_{\le}=\bigwedge_{j\ge1}([A]_j\le[B]_j),\qquad
 U_=\bigwedge_{j\ge1}([A]_j=[B]_j),                 \tag{9.1}
\]
where only finitely many padded coordinates are involved. The original
weak comparison is $U_{\le}$ together with zero-coordinate weak comparison;
the original strict comparison is $U_{\le}$ together with \emph{strict}
zero-coordinate comparison. It is neither lexicographic order nor
strictness at every coordinate.

\begin{corollary}[Whole one-ground-name source comparison]
Every finite Boolean combination of these pointwise vector equality,
weak-comparison and strict-comparison atoms is decidable under one
existential or universal quantifier over its shared positive ground
coordinate. Arbitrary original internal types and output levels are allowed.
\end{corollary}
\begin{proof}
Compute the closed guards (9.1) and put them inside their respective
pointwise Boolean atoms. Apply Corollary~\ref{cor:boolean} to the
zero-coordinate expressions. Ground free variables have precisely the
positive-natural admissible coordinate values of the original definitions
\cite[Definitions 2.5--2.8]{ww}; no extra witness variable or observer
translation is needed.
\end{proof}
One free ground name in each term is insufficient if the names differ:
the union of names across the entire query must lie in a single
$\{X:o\}$. The one-variable support hypothesis is used explicitly in
Lemma~\ref{lem:oneinput} and the range-cell rule.

For two independent inputs, a child peak depending on a small coordinate
need not dominate an affine term depending on a different large coordinate.
Our finite-zero and sign proofs therefore do not settle multivariate
weak or strict comparison. The multivariate identity theorem likewise
does not decide existence of a valuation solving a nonidentity equation;
for example $x=y$ has solutions although its normal forms differ.
Functional free variables and their admissibility constraints are outside
the source corollaries. No auxiliary numerical kernel is claimed to be
the direct assignment image of a source term.

The algorithms are mathematical decision procedures. Their finite ordinal
descent and finite syntax recursion establish termination, but evaluating
closed constants or a finite exceptional prefix can be far beyond practical
computation. The proof requires no evaluation of a giant $F_5$ value.
The role of the strong fixed control is structural: sparse peaks isolate
positive sums, and packed ordinal reserves pay entire lower profiles.

\begin{thebibliography}{99}

\bibitem{ww}Gunnar Wilken and Andreas Weiermann.
Derivation lengths classification of G\"odel's T extending Howard's assignment.
\emph{Logical Methods in Computer Science} \textbf{8}(1:19) (2012), 1--44.
\href{https://doi.org/10.2168/LMCS-8(1:19)2012}
{doi:10.2168/LMCS-8(1:19)2012}.
\href{https://lmcs.episciences.org/1073}{Official article and full text}.
The 2012 revised article is the version used throughout.

\bibitem{wwfoundations}\emph{The Wilken--Weiermann Comparison Problem: Foundations and a Guide to the Results}. Companion working manuscript, 2026.
\href{../05-wilken-weiermann-foundations/paper.pdf}{Included manuscript}.

\bibitem{wwordering}\emph{A Type-Confined Termination Ordering for G\"odel's T}. Companion working manuscript, 2026.
\href{../06-godel-t-termination-ordering/paper.pdf}{Included manuscript}.
\end{thebibliography}

\end{document}
